\documentclass[12pt]{article}
\usepackage[utf8]{inputenc}
\usepackage[T1]{fontenc}
\usepackage{hyperref}
\title{Adilabad Gondi Language Resources for K-12 Curriculum Development}
\author{Compiled by Claude Sonnet 5.5 (claude-sonnet-5-5) for Andr\'as Kornai}
\date{}
\begin{document}
\maketitle

\section{Introduction}
Adilabad Gondi (ISO 639-3 wsg, Glottocode utno1237) is a Dravidian language of Telangana, India, in the Gondi macrolanguage. Glottolog places it under Southern Gondi, with dialects including Nirmal and Rajura. The prompt suggests it may have no native speakers, but the sources found describe it as living: Omniglot reports about 300,000 speakers in 2015. It is normally written in the Telugu alphabet, and the Gunjala Gondi script, with manuscripts from the mid-18th century found in Gunjala village in Adilabad district, is also in use.

\section{Primary Educational Resources}
No Adilabad-specific textbooks or primers were verified. A Reddit post mentions a Gondi website for Adilabad with dictionaries, books and riddles in several scripts. I could not verify the site, and I saw no URL for it, so I do not list it as a resource.

\section{Linguistic and Reference Materials}
\begin{itemize}
\item Glottolog entry for wsg. \url{https://glottolog.org/resource/languoid/iso/wsg}
\item Glottolog reference for a 2015 Gondi (Adilabad) grammar sketch by Thotsam Dev Rao, Tamai Dasru, Vedma Tirupati, Arma Manik Rao and Pendur Arjun. I did not see its details. \url{https://glottolog.org/resource/reference/id/585866}
\item Omniglot, Gondi writing and script information. \url{https://omniglot.com/writing/gondi.htm}
\item The first standardized Gondi dictionary, with more than 3,000 words, was led by Shubhranshu Choudhary of CGNet Swara. It covers Gondi across regions, not Adilabad specifically. News coverage: \url{https://vikalpsangam.org/article/first-gondi-dictionary}, \url{https://www.thequint.com/amp/story/news/india/goodnews-first-ever-gondi-dictionary-tries-to-bridge-gaps} and \url{https://yourstory.com/2018/03/gondi-language-dictionary-words}.
\end{itemize}

\section{Religious and Community Resources}
ScriptureEarth lists the Gospel of Mark and the JESUS Film in Adilabad Gondi. \url{https://scriptureearth.org/00eng.php?iso=wsg}

\section{Digital and Audio Resources}
ScriptureEarth also lists Global Recordings Network audio for this language; see the URL above. No other audio, video or software resources were verified.

\section{Curriculum Development Context}
The language appears vigorous, and the prompt says it is used in education. I found no documented school programme or curriculum for it. Developers should consult the Gondi dictionary project, which is cross-regional, and the Gunjala Gondi script community. They should check the Adilabad-specific dialect before using any material.
\end{document}